%%%%%%%%%%%%%%%%
%%% exod 13 %%%
%%%%%%%%%%%%%%%%

\Chapter{13}

\Verse{1}
{
	\hP{וַיְדַבֵּר יְהֹוָה אֶל מֹשֶׁה לֵּאמֹר׃}
}
{
	\eP{And YHWH spoke to Moses, saying,}
}
{
}

\Verse{2}
{
	\hP{קַדֶּשׁ לִי כל בְּכוֹר פֶּטֶר כּל רֶחֶם בִּבְנֵי יִשְׂרָאֵל בָּאָדָם וּבַבְּהֵמָה לִי הוּא׃}
}
{
	\eP{
		“Consecrate every firstborn for me. The first birth of
		every womb of the children of Israel, of a human and of an
		animal: it is mine.”
	}
}
{
}

\Verse{3}
{
	\hJ{וַיֹּאמֶר מֹשֶׁה אֶל הָעָם זָכוֹר אֶת הַיּוֹם הַזֶּה אֲשֶׁר יְצָאתֶם מִמִּצְרַיִם מִבֵּית עֲבָדִים כִּי בְּחֹזֶק יָד הוֹצִיא יְהֹוָה אֶתְכֶם מִזֶּה וְלֹא יֵאָכֵל חָמֵץ׃}
}
{
	\eJ{
		And Moses said to the people, “Remember this day in which
		you went out from Egypt, from a house of slaves, because
		YHWH brought you out of here by strength of hand: And no
		leavened bread shall be eaten.
	}
}
{
}

\Verse{4}
{
	\hJ{הַיּוֹם אַתֶּם יֹצְאִים בְּחֹדֶשׁ הָאָבִיב׃}
}
{
	\eJ{Today you are going out, in the month of Abib.}
}
{
%	\footnote{
%		the month of Abib. This is the name of the first month of the year in
%		the Torah, the month in which spring begins. Later, the name Nisan is
%		used instead.
%	}
}

\Verse{5}
{
	\hJ{וְהָיָה כִי יְבִיאֲךָ יְהֹוָה אֶל אֶרֶץ הַכְּנַעֲנִי וְהַחִתִּי וְהָאֱמֹרִי וְהַחִוִּי וְהַיְבוּסִי אֲשֶׁר נִשְׁבַּע לַאֲבֹתֶיךָ לָתֶת לָךְ אֶרֶץ זָבַת חָלָב וּדְבָשׁ וְעָבַדְתָּ אֶת הָעֲבֹדָה הַזֹּאת בַּחֹדֶשׁ הַזֶּה׃}
}
{
	\eJ{
		And it will be, when YHWH will bring you to the land of
		the Canaanite and the Hittite and the Amorite and the
		Hivite and the Jebusite, which He swore to your fathers to
		give you, a land flowing with milk and honey, that you
		will perform this service in this month.
	}
}
{
}

\Verse{6}
{
	\hJ{שִׁבְעַת יָמִים תֹּאכַל מַצֹּת וּבַיּוֹם הַשְּׁבִיעִי חַג לַיהֹוָה׃}
}
{
	\eJ{
		You shall eat unleavened bread seven days, and on the
		seventh day is a festival to YHWH.
	}
}
{
}

\Verse{7}
{
	\hJ{מַצּוֹת יֵאָכֵל אֵת שִׁבְעַת הַיָּמִים וְלֹא יֵרָאֶה לְךָ חָמֵץ וְלֹא יֵרָאֶה לְךָ שְׂאֹר בְּכל גְּבֻלֶךָ׃}
}
{
	\eJ{
		Unleavened bread will be eaten for the seven days, and
		leavened bread shall not be seen for you, and leaven shall
		not be seen for you within all your borders.
	}
}
{
}

\Verse{8}
{
	\hJ{וְהִגַּדְתָּ לְבִנְךָ בַּיּוֹם הַהוּא לֵאמֹר בַּעֲבוּר זֶה עָשָׂה יְהֹוָה לִי בְּצֵאתִי מִמִּצְרָיִם׃}
}
{
	\eJ{
		And you shall tell your child in that day, saying,
		‘Because of that which YHWH did for me when I went out
		from Egypt.’
	}
}
{
}

\Verse{9}
{
	\hJ{וְהָיָה לְךָ לְאוֹת עַל יָדְךָ וּלְזִכָּרוֹן בֵּין עֵינֶיךָ לְמַעַן תִּהְיֶה תּוֹרַת יְהֹוָה בְּפִיךָ כִּי בְּיָד חֲזָקָה הוֹצִאֲךָ יְהֹוָה מִמִּצְרָיִם׃}
}
{
	\eJ{
		And it will become a sign on your hand and a reminder
		between your eyes for you so that YHWH’s instruction will
		be in your mouth, because YHWH brought you out from Egypt
		with a strong hand.
	}
}
{
}

\Verse{10}
{
	\hJ{וְשָׁמַרְתָּ אֶת הַחֻקָּה הַזֹּאת לְמוֹעֲדָהּ מִיָּמִים יָמִימָה׃}
}
{
	\eJ{
		And you shall observe this law at its appointed time,
		regularly.
	}
}
{
%	\footnote{
%		regularly. Literally, “from days to days,” an expression meaning “from
%		year to year.”
%	}
}

\Verse{11}
{
	\hJ{וְהָיָה כִּי יְבִאֲךָ יְהֹוָה אֶל אֶרֶץ הַכְּנַעֲנִי כַּאֲשֶׁר נִשְׁבַּע לְךָ וְלַאֲבֹתֶיךָ וּנְתָנָהּ לָךְ׃}
}
{
	\eJ{
		And it will be, when YHWH will bring you to the land of
		the Canaanite as He swore to you and to your fathers and
		will give it to you,
	}
}
{
%	\footnote{
%		He swore to you. But God did not swear it to them. Could it mean for
%		you? Probably not. Rashi, citing the Mechilta, says He swore when He
%		said, “And I’ll bring you to the land that I raised my hand …” (Exod
%		6:8). But there is no swearing of an oath to the people there. Ramban
%		suggests that the word of God itself may constitute an oath, but that is
%		not correct. The Tanak definitely distinguishes the cases in which a
%		divine oath is sworn. Ramban also takes the earlier words “He swore to
%		your fathers to give to you” (13:5) to constitute an oath to the people.
%		If an oath to the ancestors is in fact the explanation of this, then I
%		believe we should focus on the words that God says to Abraham in the
%		middle of their covenant ceremony. There the fire that expresses the
%		presence of God passes between the parts of animals, which constitutes
%		an oath ceremony; and the stated purpose of this ceremony is to confirm
%		that Abraham will possess the land. God states that after a period of
%		slavery “a fourth generation will come back here” (Gen 15:7–16). Or,
%		more generally, we may say that the promise to Abraham is repeated to
%		his descendants, the later patriarchs, showing that the oath carries
%		through the lineage to the descendants. Since the whole people of Israel
%		is a descendant of Abraham, the oath may be regarded as made to them as
%		well. I am not certain of the answer to this problem. I point it out in
%		order to make it known and in the hope that someone else may solve it.
%	}
}

\Verse{12}
{
	\hJ{וְהַעֲבַרְתָּ כל פֶּטֶר רֶחֶם לַיהֹוָה וְכל פֶּטֶר שֶׁגֶר בְּהֵמָה אֲשֶׁר יִהְיֶה לְךָ הַזְּכָרִים לַיהֹוָה׃}
}
{
	\eJ{
		that you will pass every first birth of a womb to YHWH;
		and every first birth, offspring of an animal, that you
		will have—the males—is YHWH’s.
	}
}
{
}

\Verse{13}
{
	\hJ{וְכל פֶּטֶר חֲמֹר תִּפְדֶּה בְשֶׂה וְאִם לֹא תִפְדֶּה וַעֲרַפְתּוֹ וְכֹל בְּכוֹר אָדָם בְּבָנֶיךָ תִּפְדֶּה׃}
}
{
	\eJ{
		And you shall redeem every first birth of an ass with a
		lamb, and if you will not redeem, then you shall break its
		neck. And you shall redeem every human firstborn among
		your sons.
	}
}
{
}

\Verse{14}
{
	\hJ{וְהָיָה כִּי יִשְׁאָלְךָ בִנְךָ מָחָר לֵאמֹר מַה זֹּאת וְאָמַרְתָּ אֵלָיו בְּחֹזֶק יָד הוֹצִיאָנוּ יְהֹוָה מִמִּצְרַיִם מִבֵּית עֲבָדִים׃}
}
{
	\eJ{
		And it will be, when your child will ask you tomorrow,
		saying, ‘What is this?’ that you’ll say to him, ‘With
		strength of hand YHWH brought us out from Egypt, from a
		house of slaves.
	}
}
{
}

\Verse{15}
{
	\hJ{וַיְהִי כִּי הִקְשָׁה פַרְעֹה לְשַׁלְּחֵנוּ וַיַּהֲרֹג יְהֹוָה כּל בְּכוֹר בְּאֶרֶץ מִצְרַיִם מִבְּכֹר אָדָם וְעַד בְּכוֹר בְּהֵמָה עַל כֵּן אֲנִי זֹבֵחַ לַיהֹוָה כּל פֶּטֶר רֶחֶם הַזְּכָרִים וְכל בְּכוֹר בָּנַי אֶפְדֶּה׃}
}
{
	\eJ{
		And it was when Pharaoh hardened against letting us go,
		and YHWH killed every firstborn in the land of Egypt, from
		firstborn of a human to firstborn of an animal. On account
		of this I am sacrificing to YHWH every first birth of a
		womb—the males—and I shall redeem every firstborn of my
		sons.’
	}
}
{
}

\Verse{16}
{
	\hJ{וְהָיָה לְאוֹת עַל יָדְכָה וּלְטוֹטָפֹת בֵּין עֵינֶיךָ כִּי בְּחֹזֶק יָד הוֹצִיאָנוּ יְהֹוָה מִמִּצְרָיִם׃}
}
{
	\eJ{
		And it will become a sign on your hand and bands between
		your eyes, because with strength of hand YHWH brought us
		out from Egypt.”
	}
}
{
%	\footnote{
%		a sign on your hand and bands between your eyes. The context of this and
%		the similar expression in v. 9 indicates clearly that this is a
%		metaphor, meaning that the Passover observance is to become a vivid,
%		conscious concern of the people of Israel. The occurrence of the similar
%		expression in Deut 6:8 came to be taken literally, referring to the
%		practice of wrapping one’s hand and forehead in phylacteries
%		(tephillin).
%	}
}

\Verse{17}
{
	\hRJE{וַיְהִי בְּשַׁלַּח פַּרְעֹה אֶת הָעָם וְלֹא נָחָם אֱלֹהִים דֶּרֶךְ אֶרֶץ פְּלִשְׁתִּים כִּי קָרוֹב הוּא כִּי אָמַר אֱלֹהִים פֶּן יִנָּחֵם הָעָם בִּרְאֹתָם מִלְחָמָה וְשָׁבוּ מִצְרָיְמָה׃}
}
{
	\eRJE{
		And it was, when Pharaoh let the people go,
		that {\God} did not lead them by way of the Philistines’
		land—because it was close—because {\God} said, “In case the
		people will be dissuaded when they see war, and they’ll go
		back to Egypt.”
	}
}
{
%	\footnote{
%		Philistines. Philistines are present along the coast by the time of
%		Ramesses III, who mentions them in “The War against the Sea Peoples” (c.
%		1188 B.C.E.).
%	}
}

\Verse{18}
{
	\hRJE{וַיַּסֵּב אֱלֹהִים אֶת הָעָם דֶּרֶךְ הַמִּדְבָּר יַם סוּף וַחֲמֻשִׁים עָלוּ בְנֵי יִשְׂרָאֵל מֵאֶרֶץ מִצְרָיִם׃}
}
{
	\eRJE{
		And {\God} turned the people by way of the wilderness of the
		Red Sea. And the children of Israel went up armed from the
		land of Egypt.
	}
}
{
%	\footnote{
%		by way of the wilderness of the Red Sea. They are not sent through
%		Philistine territory because they might be afraid and turn back. But
%		then they are sent to the Red Sea! Is that better? And, in fact, at the
%		Red Sea they will be afraid and wish they were back in Egypt!
%		(14:10–12). The difference is psychological. They have been slaves for
%		generations. Decision making, willpower, and responsibility have not
%		been part of their lives. One does not acquire responsibility instantly
%		when one becomes free. Facing opponents and having to take arms and
%		fight would terrify them, and their natural reaction would be to turn
%		and flee. But the Red Sea situation is different. With the Egyptian army
%		behind them and an opened sea in front of them, they do not need to make
%		a decision, exert willpower, or take responsibility. They have no choice
%		but to go forward. Footnote 13:18. Red Sea. Recent commentators and
%		translators have called this the “Sea of Reeds,” a different body of
%		water from the Red Sea; but there is no such body of water. They say
%		this because (1) the sea is called yam sûp in the Hebrew, and sûp
%		elsewhere means a reed (Exod 2:3); (2) reeds do not grow by the Red Sea;
%		and (3) they conceive of a smaller body of water than the Red Sea being
%		subject to some sort of drying or splitting, as opposed to the larger,
%		deeper Red Sea. But none of these considerations outweighs the simple
%		fact that the Tanak refers to the eastern arm of the sea—the body of
%		water now known as the Gulf of Eilat or the Gulf of Aqaba—as yam sûp as
%		well (Num 14:25; 21:4; Deut 1:40; 2:1; Judg 11:16; 1 Kings 9:26; Jer
%		49:21). There is therefore no doubt that the body of water that is
%		pictured in the story of the splitting of the sea in the Torah is the
%		body of water that is known in English as the Red Sea. I have returned
%		to this traditional English translation so it will be clear what body of
%		water is meant. I have no objection to translations that use “Reed Sea”
%		or “Sea of Suph” as long as it is clear that it is this body of water.
%		(As for its Hebrew name, yam sûp, we just do not know how that came to
%		be its name.) Footnote 13:18. armed. Or prepared to fight, or resolute
%		(Propp). But the text has just said that God does not want them to have
%		to face a fight with the Philistines. So why note that they are armed at
%		this point? Answer: (1) because they will fight the Amalekites a short
%		time later (see the comment on 17:9); and (2) in recognition of the fact
%		that one day they will have to fight for the promised land.
%	}
}

\Verse{19}
{
	\hRJE{וַיִּקַּח מֹשֶׁה אֶת עַצְמוֹת יוֹסֵף עִמּוֹ כִּי הַשְׁבֵּעַ הִשְׁבִּיעַ אֶת בְּנֵי יִשְׂרָאֵל לֵאמֹר פָּקֹד יִפְקֹד אֱלֹהִים אֶתְכֶם וְהַעֲלִיתֶם אֶת עַצְמֹתַי מִזֶּה אִתְּכֶם׃}
}
{
	\eRJE{
		And Moses took Joseph’s bones with him, because he had had
		the children of Israel swear, saying, “{\God} will take
		account of you, and you’ll bring up my bones from here
		with you.”
	}
}
{
}

\Verse{20}
{
	\hRJE{וַיִּסְעוּ מִסֻּכֹּת וַיַּחֲנוּ בְאֵתָם בִּקְצֵה הַמִּדְבָּר׃}
}
{
	\eRJE{
		And they traveled from Succoth and camped in Etham at the
		edge of the wilderness.
	}
}
{
}

\Verse{21}
{
	\hRJE{וַיהֹוָה הֹלֵךְ לִפְנֵיהֶם יוֹמָם בְּעַמּוּד עָנָן לַנְחֹתָם הַדֶּרֶךְ וְלַיְלָה בְּעַמּוּד אֵשׁ לְהָאִיר לָהֶם לָלֶכֶת יוֹמָם וָלָיְלָה׃}
}
{
	\eRJE{
		And YHWH was going in front of them by day in a column of
		cloud to show them the way, and by night in a column of
		fire to shed light for them, so as to go by day and by
		night.
	}
}
{
%	\footnote{
%		a column of cloud and a column of fire. The cloud and fire have a
%		two-sided function, both conveying God’s presence and reflecting God’s
%		mysteriousness. The narrative reports that YHWH Himself goes in front of
%		the Israelites in the wilderness, but they do not see the deity because
%		He travels in the daytime in the column of cloud and by night in the
%		column of fire. Exodus must be read with a consciousness that the cloud
%		or the fire is present at all times, a constant miracle, an introduction
%		of the cosmic in history. One must have a sense of the awe involved in
%		this phenomenon, particularly expressed in fire. It is a miraculously
%		burning bush that first draws Moses. Now the people experience a similar
%		phenomenon: a continuous, miraculous fire.
%	}
}

\Verse{22}
{
	\hRJE{לֹא יָמִישׁ עַמּוּד הֶעָנָן יוֹמָם וְעַמּוּד הָאֵשׁ לָיְלָה לִפְנֵי הָעָם׃}
}
{
	\eRJE{
		The column of cloud by day and the column of fire by night
		did not depart in front of the people.
	}
}
{
}
